% ============================================================================
% Fichier  : partie2/P02_C08_collecter_preuve_nek.tex
% Titre    : Collecter la Preuve : l'Instrument OUI/NON/NEK
% Auteur   : MINKA MI NGUIDJOÏ Thierry Emmanuel
% Version  : 1.0 -- octobre 2026
% ============================================================================
% Sources : article A0 (proposition de nécessité d'une modalité d'ignorance) ;
% article A1 (HAL hal-05701270) : espace de réponse {OUI, NON, NEK},
% métriques par item H_q, U_q, mu_q, indice de réduction du biais
% rho_q = mu^B - mu^T ; INTOSAI GUID 5100 (lettre d'extraction) ; types de
% tests et plan de tests de l'outil O7 du cours Audit_SI de l'auteur.
% Tailles d'échantillon : approximation de Poisson, zéro écart attendu.
% ============================================================================

\chapter[Collecter la Preuve]{Collecter la Preuve : l'Instrument OUI/NON/NEK}
\label{chap:preuve}
\chaptermark{Collecter la preuve}

\epigraph{%
    « Ce qui n'est pas dans le dossier\\
    n'a pas été audité. »%
}{Règle des papiers de travail}

\niveauChapitre{Intermédiaire}{%
    Maîtrise du Chapitre~\ref{chap:decomposition} (gabarit canonique, fiche
    d'assertion)%
}

\begin{boxobjectifs}
\begin{itemize}
    \item Justifier l'instrument de réponse \OUI{}~/ \NON{}~/ \NEK{} et le
          distinguer de la réponse « non applicable ».
    \item Mesurer le biais du choix forcé d'une grille binaire.
    \item Classer les preuves selon leur source et leur fiabilité.
    \item Conduire un entretien d'audit à partir des fiches d'assertion.
    \item Concevoir un plan de tests et dimensionner un échantillon.
    \item Garantir l'intégrité des extractions et tenir des papiers de
          travail jusqu'à la trace.
\end{itemize}
\end{boxobjectifs}

\bigskip

% ============================================================================
\section*{Question directrice : comment savoir qu'on sait~?}
% ============================================================================

La décomposition du domaine~D est prête (Chapitre~\ref{chap:decomposition}).
Pour chaque contrôle, une question attend sa réponse. M.~BEYALA pose à son
équipe une règle simple~: une réponse ne vaut que par la preuve qui
l'appuie, et l'on a le droit, et même le devoir, de répondre « je ne sais
pas » quand la preuve manque.

Ce chapitre organise la collecte. Il présente l'instrument de réponse
(sections~\ref{sec:c08-instrument} et~\ref{sec:c08-biais}), les preuves et
leur fiabilité (section~\ref{sec:c08-preuves}), les entretiens
(section~\ref{sec:c08-entretiens}), les tests et l'échantillonnage
(section~\ref{sec:c08-tests}), puis les extractions et les papiers de
travail (section~\ref{sec:c08-papiers}).

\separateur

% ============================================================================
\section{L'Instrument OUI / NON / NEK}
\label{sec:c08-instrument}
% ============================================================================

\subsection{Trois réponses}

L'auditologie établit, à partir de ses axiomes d'observabilité partielle, de
cognition limitée, de mesurabilité et d'interaction stratégique, que tout
instrument de réponse d'audit valide doit comporter une modalité explicite
de connaissance insuffisante~\cite{minka2026auditologie}. L'article A1 en
construit l'instrument et en étudie les propriétés~\cite{minka2026nek}.

\begin{boxdef}
\termecle{Espace de réponse trichotomique}\index{espace de réponse
trichotomique}~: ensemble $\Omega = \{\OUI, \NON, \NEK\}$ des réponses
admises à une question d'audit.
\begin{itemize}
    \item \OUI{}~: le critère est satisfait, et une preuve l'établit.
    \item \NON{}~: le critère n'est pas satisfait, et une preuve l'établit.
    \item \NEK{} (\textit{Not Enough Knowledge})~: l'information disponible
          ne permet pas de conclure, malgré des diligences normales.
\end{itemize}
\end{boxdef}

\subsection{Deux usages}

L'instrument sert à deux moments de la mission.
\begin{enumerate}
    \item \textbf{Auprès des répondants}~: une même question est posée à
          plusieurs interlocuteurs (entretiens, questionnaires). Pour chaque
          question~$q$, on obtient une distribution de réponses, dont l'article
          A1 tire trois mesures~: $\mu_q$, la part des \OUI{}~; $U_q$, la part
          des \NEK{}~; $H_q$, l'entropie~\cite{minka2026nek}.
    \item \textbf{Par l'auditeur}~: après ses tests, l'auditeur cote chaque
          contrôle \OUI{}, \NON{} ou \NEK{}, preuve à l'appui. C'est la cote
          qui entre dans le rapport.
\end{enumerate}

Les deux usages se répondent~: une question où les répondants divergent
(entropie élevée) appelle des tests~; les tests permettent ensuite à
l'auditeur de coter avec preuve.

\subsection{NEK n'est pas « non applicable »}

\begin{boxalert}
\textbf{Violation typique.} Coter \NEK{} un contrôle qui ne s'applique pas au
périmètre (par exemple, la double authentification d'un service qui n'est
pas exposé à Internet).

\smallskip
\textbf{Règle~:} un contrôle sans objet dans le périmètre est \emph{non
applicable}~: il sort de la décomposition, avec la raison écrite. \NEK{} est
réservé à un contrôle applicable dont on ne peut pas établir l'état. Les
confondre fausse le taux d'incertitude~$U$.
\end{boxalert}

% ============================================================================
\section{Le Biais du Choix Forcé}
\label{sec:c08-biais}
% ============================================================================

Une grille binaire force le répondant qui ne sait pas à choisir. Il choisit
le plus souvent la réponse attendue, ou la plus confortable~: le \OUI{}. La
grille binaire \emph{surestime} donc la conformité. L'article A1 mesure cette
surestimation par un indice de réduction du biais~\cite{minka2026nek}~:
\[
    \rho_q = \hat{\mu}_q^{B} - \mu_q^{T},
\]
où $\hat{\mu}_q^{B}$ est la part des \OUI{} obtenue avec l'instrument binaire et
$\mu_q^{T}$ celle obtenue avec l'instrument trichotomique, sur les mêmes
questions.

\begin{boxlecture}{Domaine D, binaire puis trichotomique}
Lors de la première mission, les quatre interlocuteurs avaient répondu aux
cinq questions du domaine~D avec une grille binaire~: 16~\OUI{} sur 20
réponses. Au second passage, avec \NEK{}~: 13~\OUI{}, 2~\NON{}, 5~\NEK{}.

\medskip
\textbf{Constatation.} $\hat{\mu}^{B} = 16/20 = 0{,}80$~; $\mu^{T} = 13/20 =
0{,}65$~; $\rho = 0{,}15$.

\medskip
\textbf{Analyse} (confiance moyenne). La grille binaire surestimait la
conformité du domaine de quinze points~: trois répondants qui ne savaient
pas avaient répondu \OUI{}. C'est ainsi que le premier rapport a pu conclure
à un domaine «~globalement maîtrisé~».

\medskip
\hfill\lecture{0{,}65}{0{,}25}{1{,}24}
\end{boxlecture}

% ============================================================================
\section{Les Preuves et leur Fiabilité}
\label{sec:c08-preuves}
% ============================================================================

\subsection{Suffisante et appropriée}

\begin{boxdef}
\termecle{Preuve d'audit}\index{preuve d'audit}~: information sur laquelle
l'auditeur fonde sa cotation et ses conclusions. Elle doit être
\emph{suffisante} (en quantité, au regard du risque) et \emph{appropriée}
(pertinente pour le critère, et fiable).
\end{boxdef}

\subsection{Une hiérarchie de fiabilité}

\begin{tableaucours}{cL{4.2cm}L{4.4cm}L{2.8cm}}{%
    \tableauHeader{Rang} &
    \tableauHeader{Source} &
    \tableauHeader{Exemple} &
    \tableauHeader{Fiabilité}
}
1 & Trace produite par le système, extraite par l'auditeur & Journal des
désactivations extrait en sa présence & La plus élevée \\
2 & Document externe & Relevé bancaire, attestation d'un prestataire &
Élevée \\
3 & Document interne approuvé & Compte rendu de revue signé & Moyenne \\
4 & Observation directe & L'auditeur voit le magasinier saisir une
réception & Ponctuelle \\
5 & Déclaration orale & « Les comptes sont revus chaque trimestre » & La
plus faible \\
\end{tableaucours}
\vspace{-0.8em}
\begin{center}{\small\itshape Tableau~8.1. Fiabilité des preuves selon leur
source}\end{center}

Une déclaration orale ne suffit jamais à coter \OUI{}~: elle oriente la
recherche de la preuve. Lorsque l'auditeur ne dispose que d'une
déclaration, la cote est \NEK{}, et la question devient «~où est la
preuve~?~».

% ============================================================================
\section{Les Entretiens}
\label{sec:c08-entretiens}
% ============================================================================

Le guide d'entretien (outil~O5) dérive des fiches d'assertion~: chaque
question de fiche devient une question d'entretien, classée par
interlocuteur. L'entretien d'audit n'est pas une conversation~; il obéit à
des règles.

\begin{boxmetier}{Auditeur SI}
\textbf{Conduire un entretien d'audit}
\begin{enumerate}
    \item Annoncer l'objet, la durée, et que la réponse \NEK{} est admise.
    \item Commencer par des questions ouvertes («~Comment se passe le départ
          d'un agent~?~»), puis passer aux questions de la fiche.
    \item Ne pas suggérer la réponse attendue~: «~Les comptes sont-ils
          désactivés sous 48~heures~?~» plutôt que «~Vous désactivez bien
          les comptes, n'est-ce pas~?~».
    \item Demander, pour chaque réponse affirmative, la pièce qui l'établit.
    \item Mener l'entretien à deux, l'un interrogeant, l'autre notant.
    \item Rédiger le compte rendu le jour même et le faire valider par
          l'interlocuteur.
\end{enumerate}
\end{boxmetier}

\begin{boiteRemarque}{Le répondant peut-il se cacher derrière NEK~?}
Oui~: le Chapitre~\ref{chap:fraude} l'a montré, le taux de \NEK{} peut
augmenter avec l'enjeu. C'est pourquoi l'auditeur note le contexte de
chaque réponse (entretien formel, questionnaire, auto-évaluation) et
compare les contextes. Un \NEK{} n'est jamais une fin~: il appelle un test
qui cherche la preuve directement dans le système.
\end{boiteRemarque}

% ============================================================================
\section{Tester et Échantillonner}
\label{sec:c08-tests}
% ============================================================================

\subsection{Les types de tests}

\begin{tableaucours}{cL{3.0cm}L{4.4cm}L{3.6cm}}{%
    \tableauHeader{Code} &
    \tableauHeader{Type} &
    \tableauHeader{Principe} &
    \tableauHeader{Exemple (domaine D)}
}
DOC & Revue documentaire & Lire une procédure, un compte rendu, un contrat
& Lire la procédure de départ \\
ENT & Entretien & Recueillir une déclaration, à confirmer & Interroger le
responsable RH \\
OBS & Observation & Voir le contrôle s'exécuter & Assister à une revue des
comptes \\
TEST & Test technique & Ré-exécuter le contrôle ou analyser les données &
Rapprocher départs et comptes actifs \\
MIX & Combinaison & Plusieurs types sur un même contrôle & Revue, puis
test \\
\end{tableaucours}
\vspace{-0.8em}
\begin{center}{\small\itshape Tableau~8.2. Les types de tests du plan de
tests (outil O7)}\end{center}

Un test peut viser la \emph{conception} d'un contrôle (test de cheminement
sur un cas), son \emph{efficacité} (échantillon sur la période) ou
directement les \emph{données} (tests substantifs, Chapitre~\ref{chap:audit-applications}).

\subsection{Échantillonner}

On ne teste pas toutes les exécutions d'un contrôle~: on en tire un
\termecle{échantillon}\index{échantillon}. Trois méthodes coexistent.

\begin{itemize}
    \item \textbf{Par jugement}~: l'auditeur choisit les cas les plus risqués
          (les plus gros montants, les dernières semaines de la période). Il
          ne peut pas extrapoler, mais il cible.
    \item \textbf{Aléatoire}~: chaque élément a la même chance d'être tiré~;
          l'auditeur conserve la méthode et la graine du tirage pour pouvoir
          le refaire.
    \item \textbf{Statistique par attributs}~: on fixe un niveau de confiance
          et un taux d'écart tolérable, et l'on en déduit la taille de
          l'échantillon.
\end{itemize}

Lorsque l'auditeur s'attend à ne trouver aucun écart, une approximation
courante donne la taille d'échantillon~:
\[
    n \approx \frac{F}{t},
\]
où $t$ est le taux d'écart tolérable et $F$ un facteur qui dépend du niveau
de confiance ($F \approx 2{,}3$ à 90~\%, $F \approx 3{,}0$ à 95~\%).

\begin{tableaucours}{L{4.0cm}cc}{%
    \tableauHeader{Taux d'écart tolérable} &
    \tableauHeader{Confiance 90~\%} &
    \tableauHeader{Confiance 95~\%}
}
10~\% & 23 & 30 \\
5~\% & 46 & 60 \\
\end{tableaucours}
\vspace{-0.8em}
\begin{center}{\small\itshape Tableau~8.3. Tailles d'échantillon, aucun écart
attendu (approximation)}\end{center}

Si l'échantillon révèle un seul écart, la conclusion ne tient plus~:
l'auditeur doit élargir l'échantillon ou conclure que le contrôle n'est pas
efficace. Lorsque les données sont extractibles, il vaut souvent mieux
tester l'ensemble de la population, ce qui supprime le risque
d'échantillonnage (Chapitre~\ref{chap:audit-continu}).

\begin{boxscenario}{Opération LIZBIZ. Tester D.2.1}
\textbf{Contrôle.} Désactivation des comptes des agents partis sous
48~heures (fiche du prologue).

\medskip
\textbf{Choix.} Les départs sont peu nombreux (31 sur la période)~: AUDIX
CONSEIL teste la population entière plutôt qu'un échantillon, par un
rapprochement entre la liste des départs des ressources humaines et le
journal des désactivations.

\medskip
\textbf{Constatation.} 31~départs~; 19 comptes désactivés sous 48~heures,
7 désactivés plus tard, 5 encore actifs.

\medskip
\textbf{Cotation.} \NON{}~: le critère n'est pas satisfait, et la preuve est
établie sur la population entière.
\end{boxscenario}

% ============================================================================
\section{Extractions et Papiers de Travail}
\label{sec:c08-papiers}
% ============================================================================

\subsection{Garantir l'intégrité des extractions}

Une extraction de données remise par l'audité n'est une preuve que si l'on
peut établir d'où elle vient et qu'elle n'a pas été modifiée. Le GUID~5100
de l'INTOSAI recommande qu'elle soit accompagnée d'une lettre qui en
précise la source, l'horodatage, l'empreinte numérique, les paramètres
d'extraction et l'environnement d'origine~\cite{intosai2019guid5100}.

\begin{boxoutils}{Calculer et vérifier une empreinte}
\textbf{Calculer l'empreinte d'une extraction} (Linux, ou terminal du
NDAP SUKULU)~:

\texttt{sha256sum comptes\_actifs\_2026-04-30.csv > comptes.sha256}

\medskip
\textbf{Vérifier plus tard qu'elle n'a pas changé}~:

\texttt{sha256sum -c comptes.sha256}

\medskip
Résultat attendu~: \texttt{comptes\_actifs\_2026-04-30.csv: OK}. Toute
modification du fichier, même d'un caractère, change l'empreinte.
\end{boxoutils}

\begin{boxscenario}{Opération MEFOU. La lettre d'extraction}
\textbf{Faits.} La Chambre des comptes reçoit de l'ANPP l'extraction des
comptes du système de paiement.

\medskip
\textbf{Contenu de la lettre} signée par la cheffe de la division
informatique~: base de production, requête utilisée, date et heure
d'extraction, nombre de lignes, empreinte SHA-256 du fichier.

\medskip
\textbf{Analyse} (confiance élevée). Le magistrat peut prouver que les écarts
qu'il relèvera proviennent bien des données de l'agence, et non d'un fichier
retouché~: la constatation devient opposable.
\end{boxscenario}

\subsection{Les papiers de travail}

\begin{boxdef}
\termecle{Papiers de travail}\index{papiers de travail}~: ensemble des
documents établis ou recueillis par l'auditeur, qui retracent ses
travaux, ses preuves et le raisonnement qui conduit à ses conclusions.
\end{boxdef}

\begin{boxprocedure}
\textbf{Le papier de travail d'un test.} Pour qu'un tiers puisse refaire le
test et aboutir à la même cotation, chaque papier de travail indique~:
\begin{itemize}
    \item sa référence, et celle du contrôle dans la décomposition (par
          exemple PT-14, contrôle D.2.1)~;
    \item l'objectif du test, le critère et la population ou l'échantillon,
          avec sa méthode de sélection~;
    \item les preuves examinées, avec leur source et leur empreinte, et les
          traces correspondantes~;
    \item les résultats, les écarts relevés et la cotation \OUI{}, \NON{} ou
          \NEK{}, justifiée~;
    \item le nom de l'auditeur, la date, et le visa du réviseur.
\end{itemize}
\end{boxprocedure}

\begin{boxmetier}{Commissaire aux comptes}
\textbf{S'appuyer sur les travaux de l'auditeur SI}

Le commissaire aux comptes de LiZbiZ peut utiliser les travaux d'AUDIX
CONSEIL sur le contrôle d'accès pour fonder sa propre opinion sur les
comptes. Il ne le fera que si les papiers de travail lui permettent
d'apprécier la qualité des preuves~: méthode d'échantillonnage, intégrité des
extractions, cotations justifiées. Un \NEK{} documenté lui est plus utile
qu'un \OUI{} sans preuve.
\end{boxmetier}

\separateur

\begin{boxresume}
\textbf{Ce qu'il faut retenir de ce chapitre~:}
\begin{itemize}
    \item \textbf{Trois réponses}~: \OUI{} et \NON{} s'appuient sur une
          preuve~; \NEK{} déclare une connaissance insuffisante~; « non
          applicable » sort du périmètre et n'est pas \NEK{}.
    \item \textbf{Le choix forcé surestime la conformité}~: l'indice
          $\rho = \hat{\mu}^{B} - \mu^{T}$ mesure ce biais.
    \item \textbf{Les preuves n'ont pas la même fiabilité}~: la trace extraite
          par l'auditeur l'emporte sur la déclaration orale, qui ne suffit
          jamais à coter \OUI{}.
    \item \textbf{L'entretien} suit les fiches d'assertion, sans suggérer la
          réponse, et demande la pièce.
    \item \textbf{Échantillonner} se fait par jugement, au hasard ou par
          attributs~; quand c'est possible, on teste la population entière.
    \item \textbf{L'intégrité des extractions} se prouve par leur empreinte~;
          les papiers de travail permettent à un tiers de refaire chaque test.
\end{itemize}
\end{boxresume}

\bigskip

\begin{boxexo}[title={\ding{45}\quad Exercice~8.1. OUI, NON, NEK ou non applicable~? \niveaubase}]
Cotez chaque situation et justifiez.
\begin{enumerate}[(a)]
    \item Le responsable informatique affirme que les journaux sont revus
          chaque semaine, sans pouvoir le montrer.
    \item Le journal des désactivations montre que tous les départs du
          trimestre ont été traités sous 48~heures.
    \item Le contrôle de la double authentification à distance, alors
          qu'aucun accès à distance n'existe.
    \item Le journal de revue des anomalies existe, mais il est vide pour
          les trois derniers mois, alors que la procédure exige une revue
          hebdomadaire.
\end{enumerate}
\end{boxexo}

\begin{boxexo}[title={\ding{45}\quad Exercice~8.2. Dimensionner et mesurer \niveaumoyen}]
\begin{enumerate}[(a)]
    \item Quelle taille d'échantillon retenir pour tester, avec une confiance
          de 95~\% et un taux d'écart tolérable de 5~\%, la validation des
          ajustements de stock de LiZbiZ~?
    \item Le test révèle un écart. Que doit faire l'auditeur~?
    \item Sur une question posée à dix répondants, la grille binaire donnait
          8~\OUI{}~; la grille trichotomique donne 5~\OUI{}, 1~\NON{},
          4~\NEK{}. Calculez $\rho$, $U$ et $H$, et interprétez.
\end{enumerate}
\end{boxexo}

\begin{boxexo}[title={\ding{45}\quad Exercice~8.3. Le plan de tests du sous-domaine D.4 \niveauexpert}]
\begin{enumerate}[(a)]
    \item À partir de la fiche d'assertion D.4.2 (Chapitre~\ref{chap:decomposition}),
          rédigez trois tests de types différents.
    \item Pour chacun, précisez la preuve attendue, sa fiabilité et la trace
          recherchée.
    \item Rédigez trois questions du guide d'entretien destinées au
          responsable de la sécurité.
    \item Rédigez la lettre d'extraction que vous demanderez pour la liste
          des comptes d'administration.
\end{enumerate}
\end{boxexo}

\bigskip

\begin{boxinfo}
\textbf{Transition.} Les réponses et les preuves sont collectées. Il faut
maintenant conclure~: calculer les mesures de chaque domaine, décider ce
qu'elles permettent d'affirmer, et rédiger les observations. C'est l'objet du
Chapitre~\ref{chap:conclure}.
\end{boxinfo}
